\section{Appendix Overview}

The appendix contains the complete pre-results experimental contract.
It is intentionally detailed so that model-specific implementation choices cannot be changed silently after test outcomes are observed.

\section{Trajectory Telemetry Contract}

\paragraph{Eligibility.}
At every observation, predictive statistics are pooled only over positions that are active for the current block and unresolved.
Prompt tokens, inactive future blocks, and committed tokens are excluded.

\paragraph{Observation time.}
Telemetry is captured after the model produces logits and before the sampler commits/remasks tokens.
This makes every uncertainty feature pre-decision with respect to that denoising step.

\paragraph{JS drift.}
Exact Jensen--Shannon drift is calculated from consecutive full-vocabulary distributions only when a position is unresolved in both observations.
One previous probability tensor is retained; no full trajectory of vocabulary distributions is stored.

\paragraph{Decision-time features.}
For each generated position, decision-time entropy, confidence, and margin use its last unresolved observation.
The baseline uses mean decision-time entropy.

\paragraph{Commitment time.}
If a token first appears committed at observation $k$, the decision is assigned to observation $k-1$.
A token still unresolved in the final pre-transfer observation is assigned the final active decision index.

\section{Frozen Feature Definitions}

\begin{table*}[t]
\centering
\small
\begin{tabular}{lll}
\toprule
Family & Feature & Definition\\
\midrule
Decision & decision\_entropy\_mean & mean normalized entropy immediately before each token commitment\\
Distribution & entropy\_mean & mean normalized entropy over active unresolved positions\\
Distribution & entropy\_p90 & 90th percentile of normalized entropy\\
Distribution & entropy\_slope & linear early-to-late trend of step-level entropy\\
Distribution & margin\_mean / margin\_min & top-1 minus top-2 probability margin\\
Stability & js\_mean / js\_p90 & consecutive-step normalized JS drift\\
Stability & js\_slope / late\_js\_mean & temporal and late-stage JS behavior\\
Stability & revision\_rate & top-1 proposal changes while unresolved\\
Stability & late\_revision\_rate & revision rate in the final quarter of observations\\
Persistence & mask\_auc & unresolved fraction integrated over active observations\\
Persistence & late\_mask\_fraction & unresolved fraction in late observations\\
Persistence & commitment\_time\_mean & normalized mean decision index\\
\bottomrule
\end{tabular}
\caption{Frozen primary feature dictionary.}
\end{table*}

\section{Calibration and Leakage Firewall}

Generation and scoring are two separate stages.
The raw JSONL record stores prompt hash, model/runtime revision, seed, generated text/token IDs, trajectory features, step summaries, and runtime metadata.
It must store neither the reference answer nor correctness.
The scorer later joins correctness only after verifying example identity and prompt hash.
Calibration and test IDs are disjoint and namespaced by dataset and immutable source split.

Feature selection, calibrator family, adaptive thresholds, and ablation definitions are frozen before final test evaluation.
The primary calibrator is standardized logistic regression.
Isotonic regression is a sensitivity analysis only when the calibration partition is large enough to support it.

\section{Model-Specific Parity Gates}

\subsection{LLaDA-MoE}

The instrumented sampler mirrors the public LLaDA low-confidence masked-diffusion loop.
A 20-prompt deterministic reference run without telemetry is compared against the instrumented run using exact generated token IDs.
Mask token ID, attention mask, transfer counts, model/tokenizer revision, and seed must match.
Primary results are blocked until exact parity passes.

\subsection{Dream}

Dream exposes a token-generation hook through its public diffusion interface.
We compare a native no-hook reference against the hook-enabled telemetry path on identical prompts, revision, seed, and decoding parameters.
Exact token-ID parity is required.

\subsection{SDAR}

SDAR operates blockwise.
Telemetry is captured from the active generated positions of each block before token transfer.
The adapter is compared against the pinned upstream block-diffusion runtime.
The aggregate feature semantics and block boundaries are recorded in run metadata.
Primary SDAR numbers are blocked until parity passes.

\subsection{FLARE / HybridDiffusion}

FLARE is used primarily for RQ4.
The native runtime trace is parsed externally rather than modifying the serving engine.
Top-$k$ trace entropy is treated only as a proxy and is never mislabeled as full-vocabulary entropy.
Same-checkpoint AR/diffusion disagreement is a diagnostic, not a correctness label.

\section{Dataset Split Policy}

\begin{table}[h]
\centering
\small
\begin{tabular}{lll}
\toprule
Dataset & Calibration & Test\\
\midrule
GSM8K & official train & official test\\
ARC-Challenge & official validation & official test\\
MMLU-Pro & official validation & official test\\
GPQA-Diamond & frozen internal holdout & frozen internal holdout\\
TruthfulQA & frozen internal holdout & frozen internal holdout\\
MATH-500 & frozen internal holdout & frozen internal holdout\\
\bottomrule
\end{tabular}
\caption{Planned split roles. Internal holdouts are created before model calls.}
\end{table}

All prepared example IDs are namespaced as
\texttt{dataset:source-split:item-id}
to prevent false overlap when official splits reuse integer row indices.

\section{Correctness Scoring}

Multiple-choice tasks use exact option-letter matching after extracting the explicit final answer.
GSM8K uses normalized numeric exact matching against the official answer field.
MATH-500 uses a pinned symbolic equivalence verifier; parser failures are conservatively counted as incorrect.
No LLM judge is used for the primary correctness labels.

\section{Primary Statistical Tests}

For each model--dataset cell, the primary nested comparison fits:
\begin{align}
\mathcal{M}_0 &: H^{\mathrm{dec}} \rightarrow P(\mathrm{error}),\\
\mathcal{M}_1 &: [H^{\mathrm{dec}}, z_{\mathrm{traj}}] \rightarrow P(\mathrm{error}).
\end{align}
Both are fit only on calibration examples and evaluated on the same test examples.
Paired class-stratified bootstrap resampling estimates confidence intervals for AUROC, AUPRC, and AURC differences.
AURC improvements are sign-normalized so positive always means the trajectory model is better.

\section{Feature-Family Ablations}

We freeze:
\begin{enumerate}
\item decision only;
\item decision + distribution features;
\item decision + stability features;
\item decision + persistence features;
\item full model;
\item full minus JS drift;
\item full minus revisions;
\item full minus persistence.
\end{enumerate}
Every ablation is calibrated on the same calibration partition and evaluated on the same test examples.

\section{Repeated-Sampling Baseline}

The expensive baseline draws five independent generations under a frozen sampling configuration.
Without consulting labels, final answers are canonicalized and summarized by:
variation ratio, normalized vote entropy, and unique-answer fraction.
These features are calibrated on the calibration partition.
All additional model calls are included in the compute report.
For MATH-500, repeated-sampling voting remains disabled until algebraic-equivalence clustering is available.

\section{Adaptive Compute}

At the base budget, trajectory features are calibrated into error probability.
To prevent optimistic threshold fitting, calibration-set policy scores are generated by stratified cross-fitting.
The threshold is chosen under a frozen maximum refinement fraction.
A final calibrator is then trained on all calibration examples and applied to test.
The threshold is not retuned on test.

We report policy accuracy, refinement fraction, mean compute, and paired bootstrap intervals for accuracy differences versus base, fully refined, and expected random refinement.

\section{Run Provenance}

Every reported number must be associated with:
\begin{itemize}
\item repository commit;
\item model and tokenizer repository plus immutable revision;
\item upstream runtime commit where applicable;
\item dataset repository/config/split and immutable revision/fingerprint;
\item prompt-template version;
\item seed and decoding configuration;
\item Python, PyTorch, Transformers, CUDA, and GPU information;
\item wall-clock time and peak allocated GPU memory;
\item raw and scored artifact SHA-256 hashes.
\end{itemize}

\section{Planned Result Tables}

The final appendix will contain:
per-model/per-dataset uncertainty metrics;
calibration sensitivity;
feature-family ablations;
denoising-budget sensitivity;
adaptive-compute results;
parity audits;
runtime/memory overhead;
and a failure taxonomy.
No cell is populated without a corresponding immutable result artifact.
